\documentclass[../main.tex]{subfiles}

\begin{document}

\begin{abstract}
Algorithms are a fundamental component of computer science education and play an important role in developing computational thinking, problem-solving
  skills, and procedural reasoning among secondary students.
As computer science programs continue to expand in secondary education, researchers have increasingly explored effective ways to teach, assess, and
  support algorithmic learning.
This literature review examines research on algorithm instruction in secondary education, focusing on instructional strategies, student learning
  challenges, assessment practices, and teacher preparation.

The literature indicates that effective algorithm instruction emphasizes conceptual understanding, abstraction, decomposition, and algorithmic
  reasoning rather than programming syntax alone.
Studies highlight the benefits of visualization tools, scaffolded learning experiences, authentic problem-solving contexts, collaboration, and guided inquiry. 
Other benefits include the identification of common student difficulties such as misconceptions, abstract reasoning challenges, and cognitive load.
Research further suggests that comprehensive assessment approaches and strong teacher preparation are essential for supporting student learning.
Overall, successful algorithm education depends on the integration of effective teaching methods, meaningful assessment practices, and well-prepared
  educators, although additional research is needed to determine the most effective approaches across diverse secondary educational settings.
\end{abstract}

% \begin{IEEEkeywords}
% Secondary Education, Computer Science Education, Literary Review
% \end{IEEEkeywords}

\end{document}
